\scanpage{203}

% B01 / text / upper edge
\mspage{104}
\deleted{\unsure{$\Pi$|$\varphi$}}

% B02 / display / top, (25)
\[
 \Pi=\Pi_R=\{\,\varphi:S\simeq S^*\mid\varphi\text{ est }R\text{-symétrique}\,\} \tag*{(25)}
\]

% B03 / text / upper, choice of gauge
des $R$-jauges. \deleted{On a \inserted{dans} dans ceci.} On voit que
( pour un tel choix $\varphi$

% B04 / display / upper, (26)
\[
 \varphi_0:S\simeq S^*\qquad\varphi_0\in\Pi\ , \tag*{(26)}
\]

% B05 / text / upper to middle, bijections
les applications bijectives $\varphi:S\simeq S^*$ correspondent
biunivoquement aux permutations $\lambda\in\unsure{\mathfrak S|\mathfrak G}_{S}$,
via

% B06 / display / upper middle, (27), permutation below lambda
\[
 \begin{array}{c@{\quad}l}\lambda &\longmapsto\ \varphi=\varphi_\lambda=\varphi_0\circ\lambda\\[-2pt]
 \rotatebox{-90}{$\in$}&\\[-4pt]\mathfrak S_S&\end{array} \tag*{(27)}
\]

% B07 / text / middle, second permutation
ou encore aux $\lambda^*\in\mathfrak S_S$

% B08 / display / middle, (27 star)
\[
 \begin{array}{c@{\quad}l}\lambda^*&\unsure{\longmapsto|\vdash}\ \varphi=\varphi_{\lambda^*}=\deleted{\varphi}\lambda^*\circ\varphi_0\\[-2pt]
 \rotatebox{-90}{$\in$}&\\[-4pt]\mathfrak S_{S^*}&\end{array} \tag*{$(27^*)$}
\]

% B09 / text / middle, introducing (28)
la relation entre $\lambda$ et $\lambda^*$ étant donnée par

% B10 / display / middle, (28), brace
\[
 \left\{\begin{aligned}
 \lambda^*&=\varphi_0(\lambda)=\varphi_0\lambda\varphi_0^{-1}\\
 \lambda&=\varphi_0^{-1}(\lambda^*)=\varphi_0^{-1}\lambda^*\varphi_0
 \end{aligned}\right. \tag*{(28)}
\]

% B11 / text / lower middle, condition
Ceci dit, si $\varphi$ correspond à $\lambda$, pour que
$\varphi$ soit une $R$-jauge, il f. et s. que l'on ait

% B12 / display / lower, (29), brace and interlinear formula
\[
 \left\{\begin{aligned}
 &\text{la relation}\quad(s,\overset{\inserted{\lambda^*\varphi_0(t)}}{\overbrace{\varphi_0(\lambda(t))}})\in R
     \quad\text{est symétrique en }s,t,\\
 &\text{i.e. la relation }\{s,\lambda(t)\}\in A_0\quad\text{symétrique en }(s,t)\\
 &\deleted{\text{est}}\ \text{elle est antiréflexive i.e.}\quad
       \forall s\in S,\ \text{on a }\{s,\lambda(s)\}\notin A_0.
 \end{aligned}\right. \tag*{(29)}
\]
% slanted
\margin{antiréflexivité pas essentielle, si on \unsure{accepte|accept.|exceptées} des graphes impropres}

% B13 / text / lower, symmetry
La symétrie s'écrit aussi sous la forme

% B14 / display / lower, (29 prime), equivalence and implication
\[
 \begin{aligned}
 &\forall(s,t)\in S\times S,\ \text{on a }\{s,t\}\in A_0
 \Longleftrightarrow\{\lambda^{-1}s,\lambda^{\deleted{\unsure{-1|1}}}t\}\in A_0\\
 &\quad(\text{il suffit que }\forall s,t\in S\times S,\quad
 \{s,t\}\in A_0\Longrightarrow\{\lambda^{-1}s,\lambda t\}\in A_0 .
 \end{aligned}\tag*{$(29')$}
\]

% B15 / text / lower left, before last display
\deleted{ou encore} ou \deleted{\unsure{encore|en}}

% B16 / display / foot, (29 double prime)
\[
 \forall s\in S,\ \text{on a }\quad
 \lambda\underbrace{(\mathrm{Et}_{A_0}(\unsure{s|t}))}_{\varphi_0(s)}
 =\underbrace{\mathrm{Et}(\lambda^{-1}s)}_{\varphi_0(\lambda^{-1}s)} \tag*{$(29'')$}
\]

% B17 / text / bottom, unfinished sentence
-- sous cette forme, est bien mis en évidence l'action

% archive sheet number: 202
